%% kapa: exam=1 round=21 year=2010 period=2 form=A subject=accounting subject_ko=회계학 seq=34 number=34
%% source: Q-Net 감정평가사 기출문제 (cst00309) artlSeq=5024181 file=extracted/1008/5024181/2010년 제21회 감정평가사 1차시험 기출문제/2교시.hwp
%% conversion: kordoc 4.24.51 -> tools/gen1.py (auto); review=proofread-gpt-sol; flags=none
\begin{question}{34}
(주)한국은 제품 A, B, C 세 가지 결합제품을 생산하고 있다. 관련 자료는 다음과 같다. 결합원가가 분리점에서의 상대적 판매가치에 의하여 배분된다면 제품 B에 배분되는 결합원가는 얼마인가? (단, 아래 표에서 세 가지 제품 모두 분리점에서의 판매가치를 알 수 있으며, 재공품은 없다고 가정함)

{\small\centering
\begin{tabular}{|l|r|r|r|r|}
\hline
 & \multicolumn{1}{c|}{제품 A} & \multicolumn{1}{c|}{제품 B} & \multicolumn{1}{c|}{제품 C} & \multicolumn{1}{c|}{합계} \\
\hline
생산수량(개) & 12,000 & 7,000 & 6,000 & 25,000 \\
\hline
결합원가(원) & ? & ? & 38,000 & 200,000 \\
\hline
\makecell[l]{분리점에서의\\판매가치(원)} & 120,000 & ? & ? & 500,000 \\
\hline
\makecell[l]{추가가공시의\\추가원가(원)} & 25,000 & 15,000 & 11,000 & 51,000 \\
\hline
\makecell[l]{추가가공후의\\판매가치(원)} & 230,000 & 200,000 & 180,000 & 610,000 \\
\hline
\end{tabular}\par}

\begin{choices}
  \item ￦114,000
  \item ￦120,000
  \item ￦124,000
  \item ￦128,000
  \item ￦130,000
\end{choices}
\end{question}
\answer{1}
